\section{Considerações Finais}

Este artigo propõe investigar como os determinantes do retorno econômico esperado e o ambiente institucional condicionam a conversão de vegetação nativa em Mato Grosso. O recorte preserva a hipótese H1 da tese e seus desdobramentos relativos à capacidade produtiva local e ao condicionamento institucional. A contribuição pretendida está em examinar esses mecanismos sob uma base territorial comparável, distinguindo atividades, biomas e destinos da conversão.

A literatura revisada oferece razões para investigar o papel dos preços e das instituições, mas também impede antecipar uma resposta única. Incentivos favoráveis à produção podem ampliar a abertura de áreas ou reorganizar usos já existentes; medidas institucionais podem conter a conversão, alterar sua composição ou deslocar sua ocorrência. A análise econométrica deverá confrontar essas possibilidades com as previsões formuladas, observando a validade das medidas e das condições de identificação.

Isto posto, as bases reunidas constituem um ponto de partida, cujo aproveitamento depende da construção das conversões brutas, da harmonização territorial e da validação temporal das exposições produtivas. A fiscalização permanece uma dimensão potencialmente endógena, e a disponibilidade de registros administrativos não deve ser confundida com identificação de sua eficácia. O período de 2001--2024 é candidato, com definição final condicionada a essas etapas.

No plano da tese, compreender como decisões econômicas alteram o estoque ambiental prepara a investigação de como esse estoque participa das condições produtivas. A sensibilidade territorial dos incentivos e a distribuição de seus efeitos ao longo do tempo orientam essa articulação. Nesta etapa, contudo, as implicações de política permanecem condicionadas aos resultados a serem obtidos, sem pressupor a falha dos instrumentos existentes ou a superioridade de um arranjo de governança.

